\documentclass[12pt, a4paper]{article}
\usepackage[T1]{fontenc}
\usepackage[utf8]{inputenc}
\usepackage{authblk}
\usepackage{ragged2e}
\usepackage{graphicx}
\usepackage[margin=1in, includefoot]{geometry}
\usepackage{amsmath}
\usepackage{setspace}
\usepackage{longtable}

\onehalfspacing

\newcommand{\vecw}{\underline{\textit{w}}}
\newcommand{\fw}{\textit{f}(\vecw)}

\DeclareRobustCommand{\ftnlogo}{%
  \begingroup\normalfont
  \includegraphics[height=4\fontcharht\font`\B]{FTN_LogoResized.png}%
  \endgroup
}

\DeclareRobustCommand{\unslogo}{%
  \begingroup\normalfont
  \includegraphics[height=4\fontcharht\font`\B]{UNS_LogoResized.png}%
  \endgroup
}

\begin{document}
    \begin{titlepage}
        \begin{center}
        \huge{\textbf{\ftnlogo Univerzitet u Novom Sadu \unslogo\\ [12pt]  Fakultet tehničkih nauka}}\\
        \vspace{55mm}
        \LARGE{\textbf{Dokumentacija za projektni zadatak}}\\
        \vspace{35mm}
        Januar 2021.
        \vspace{35mm}
        \end{center}
        \begin{flushright}
            Studenti:\hfill Ilija Kalinić, SW65/2019\\
                            Zoran Bukorac, SW40/2019\\
                            Matija Zarić, SW24/2019\\
            \vspace{3pt}
            Predmet:\hfill Nelinearno Programiranje i Evolutivni Algoritmi\\
            \vspace{3pt}
            Broj projektnog zadatka:\hfill br. 12\\
            \vspace{3pt}
            Tema projektnog zadatka:\hfill PSO algoritam, "black-box"optimizacija\\
        \end{flushright}
        
    
    \end{titlepage}

    \section{Opis problema}
    Zadata je funkcija \fw\space za koja je definisana preko veštačke neuronske mreže, pa je možemo smatrati nepoznatom.
    Poznato je da je \vecw\space vektor dimenzije 60, dakle:
    \begin{equation}
    \vecw = ( w_1, w_2, w_3, \ldots, w_{58}, w_{59}, w_{60} ),
    \end{equation}
    i da je vrednost kriterijuma realan, jednodimenzionalan broj. \\
    Funkcija \fw\space je nepoznata u analitičkom smislu (ne postoji matematička formula koja
    opisuje tu funkciju), ali možemo evaluirati njenu vrednost za proizvoljnu vrednost argumenata
    \vecw.
    U priloženoj datoteci \textit{ann\_criterion.py}\space se nalazi kod za veštačku neuronsku
    mrežu odnosno kod za izračunavanje vrednosti našeg kriterijuma optimalsnosti. \\

    Potrebno je optimizovati ulazne vrednosti koristeći PSO (Particle Swarm Optimization) algoritam
    kako bismo minimizovali vrednost zadate nepoznate funkcije, tj. kriterijuma optimalnosti \fw. 
    

    
    \begin{figure}[h]
        \centering
        \includegraphics[scale=0.36]{Blackbox3D-withGraphs.png}
        \caption{Skica "Black-boxa"\footnotemark}
        \label{fig:Black_box}
    \end{figure}
    
        \footnotetext{Autor: Krauss - Own work, CC BY-SA 4.0,
        https://commons.wikimedia.org/w/index.php?curid=45608217}
    
    \section{Uvod}
    
    \subsection{Kratak opis Particle Swarm Optimization algoritma}
    Particle Swarm Optimization algoritam je genetički algoritam koji donekle simulira socijalna ponasanja zivotinja (najvise ptica) u velikim grupama. 
    Glavni cilj algoritma PSO je da optimizuje zadat prostor pretrage koristeci roj cestica u tom prostoru. 
    Roj čestica se može još i nazvati populacijom. Svaku česticu karakteriše njen položaj u prostoru, koji je ujedno i jedno potencialno rešenje problema. 
    Inicijalne pozicije svih čestica u prostoru se određuju nasumično. 
    Čestice određuju poziciju u sledećoj iteraciji na osnovu trenutne vrednosti kriterijuma optimalnosti (inercioni faktor), sopstvene najoptimalnije vrednosti (kognitivni faktor), najoptimalnije vrednosti u celoj populaciji (socijalni faktor).
    Uticaj sopstvene najoptimalnije vrednosti treba da bude veći u početnim iteracijama (čestica sama traži optimum na osnovu svojih "najboljih iskustava"), a kako teče optimizacija da slabi, a uticaj najoptimalnije vrednosti u celoj populaciji treba da bude obrnut, manji u početku a veći na kraju optimizacije (čestica se posle pridružuje ostatku populacije u pronalaženju minimuma). \\
    \begin{figure}[h]
        \centering
        \includegraphics[scale=0.5]{PSO - skica kretanja.png}
        \caption{Skica kretanja čestice u postupku optimizacije}
        \label{fig:Walking sketch}
    \end{figure}
    
    %\subsection{}
    
    \section{Implementacija}
    Uzor za implementaciju PSO algoritma je bilo predavanje profesora Rapaić na kursu\\
    Nelinearno Programiranje i Evolutivni Algoritmi.\\
    
    Implementacija je odrađena u programskom jeziku \textit{\textbf{Python}}, uz dve ugrađene     biblioteke \textit{math}, \textit{random}, \textit{time}\space i dodatnom bibliotekom za numeričke
    operacije \textit{numpy}.\\
    U implementaciji se koriste 3 datoteke koji sadrže potrebne klase i    kod za algoritam (\textit{options.py}, \textit{particle.py} i \textit{pso.py}) , priložena
    datoteka koja sadrži kriterijum optimalnosti (\textit{ann\_criterion.py} i skripta za    pokretanje programa (\textit{main.py}).   \\
    
    \newpage
    
    \subsection{Implementacija pomoćnih struktura}
    
        Implementirane su klase:
        \begin{itemize}
            \item \textbf{MyOptions} - sadrži atribute za sve opcije i parametre koji su potrebni za izvršavanje optimizacionog algoritma.
            \item \textbf{Particle} - sadrži atribute u kojima se čuvaju vrednosti tačaka u u svim iteracijama, brzine čestice u svim iteracijama, tačka u kojoj je čestica bila najoptimalnija, kao i vrednost kriterijuma optimalnosti u toj tački i broj dimenzija prostora.\\
            Dodatno postoje i funkcije za evaluiraje trenutne tacke u kojoj se nalazi čestica, i pomoćne funkcije koje ažuriraju ostale atribute.
            \item \textbf{PSO} - sadrži glavnu petlju za izvršavanje optimizacionog algoritma, a u atributima čuva vrednosti najbolje tačke i vrednosti kriterijuma optimalnosti u toj tački, populaciju (skup svih čestica) u  listi.
        \end{itemize}
    
    \subsection{Glavni deo programa}
        Glavni fajl koji povezuje sve strukture i funkcije je skripta \textit{main.py}. U njoj se inicijalizuje objekat \textit{Options}\space koji sadrži potrebne parametre za socijalni, kognitivni i inercioni faktor, zatim korisnik određuje broj iteracija i broj populacije (upisom tih vrednosti u sam fajl), pa se skripta može pokrenuti.\\
        Na kraju izvršenja optimizacije se na konzoli ispiše optimalna tačka i njena vrednost, kao i okvirno vreme za koliko se izvršila optimizacija upotrebom modula \textit{time}.
        
    \subsection{Klasa \textit{Options}}
    Veoma je značajna klasa \textit{Options} iz razloga sto se u njoj nalaze pocetne i
    kranje vrednosti inercionog, kognitivnog i socijalnog faktora kretanja svake cestice.
    Profesor Rapaic je na svom predavanju o PSO algoritmu objasnio da se od nastanka algoritma eksperimentalnim putem doslo do optimalnih podesavanja za pocetne i krajnje vrednosti ovih parametra.
    

    
    \section{Zaključak}
    Puštanjem skripte da izvrši optimizaciju nekoliko puta, svaki put uz podešavanje broja iteracija i brojnosti populacije smo došli do najoptimalnijeg zabeleženog rešenja:\\

    \newpage
    \begin{center}
       \vecw\space \textbf{=}\space (-1.705, -1.594, -8.866, -1.112, -0.347, -1.810,\\
        -0.996, -1.206, 0.432, -1.427, 0.015, 0.856,\\
        2.411, -0.674, 0.318, 1.397, -0.698, 0.965,\\
        0, 7.238, -1.523, 0.028, -0.660, 0.996,\\
        1.197, -0.887, 0.036, -1.417, -0.735, 1.044,\\
        -1.075, 1.648, 2.054, -0.028, -0.207, -4.721,\\
        -6.096, -1.138, 0.329, -3.026, -0.700, 1.835,\\
        2.744, 0.012, 0.360, -0.873, -0.946, 0.437,\\
        -2.413, -1.176, -0.404, 0.510, -1.886, 3.665,\\
        -0.772, -1.806, 1.316, -0.803, -1.501, 1.126)\\
        \fw \textbf{=}\space 0.011\\ 
        Gde je \textbf{broj iteracija} bio \textbf{500},\\
        a \textbf{brojnost populacije 300}\\
    \end{center}
    
    Kod najvećeg broja rezultata optimizacije, dobili smo slične optimume sa razlikama
    u redu $10^{-3}$.\\
    Pored preciznosti optimizacije, dodatan problem pravi i vreme za koje će se izvršiti
    optimizacija. Na ovo utiče \underline{broj iteracija} i \underline{brojnost} populacije. Ustanovljeno je da u slučaju optimizacije zadatog kriterijuma
    optimalnosti, mnogo više utiče \textbf{brojnost populacije} nego broj iteracija u brzini
    izvršavanja optimizacije, pretpostavljamo zbog toga što je dimenzionalnost kriterijuma
    optimalnosti 60.
    
    \vspace{20mm}
    
    \begin{figure}[h]
        \centering
        \includegraphics[scale=0.5]{Options_class.png}
        \caption{Isečak koda iz fajla options.py\\cp - kognitivni faktor, cg - socijalni faktor, w - inercioni faktor}
        \label{fig:code_snippet}
    \end{figure}
    
    
\end{document}
